\documentclass[12pt]{article}

\usepackage[portuguese]{babel}
\usepackage{float}
\usepackage{listings}
\usepackage[T1]{fontenc}
\usepackage{indentfirst}
\usepackage{amsfonts, amssymb, amsmath}
\usepackage{graphicx}
\usepackage{enumitem}
\usepackage[colorlinks=true, allcolors=blue]{hyperref}
\usepackage[letterpaper,
            top=1.5cm,
            bottom=1.5cm,
            left=2.75cm,
            right=2.75cm,
            marginparwidth=1.75cm]{geometry}

\title{Probabilidade e Estatística}
\author{Prof. Douglas}
\date{\today}

\begin{document}
\maketitle

\section{Introdução a Estatística}
\begin{itemize}
    \item A estatistica envolve a coleta, organização, análise, interpretação e 
    apresentação de dados
    \item Seu proposito central é extrair conhecimento 
    \item População: Conjunto completo de elementos investigados
    \item Amostra: Subconjunto representativo da população
    \item Variável de interesse: 
\end{itemize}

\section{Análise Exploratória de Dados}
    \subsection{População e Amostra}
    \begin{itemize}
        \item População -- Conjunto completo de individuos ou unidades de 
        interesse sobre os quais queremos obter informações.
        \item Amostra -- Subconjunto selecionado da população, usado para obter 
        informações e tirar conclusões sobre ela.
        \item Unidade Amostral -- Cada elemento individual da amostra. Unidade 
        sobre a qual as observações ou medições são realizadas.
    \end{itemize}

    \subsection{Amostragem}
    \begin{itemize}
        \item Probabilística
        \begin{itemize}
            \item Envolve seleção aleatória, garantindo que cada elemento da 
            população tenha uma chance conhecida de ser incluído na amostra
            \item Considerado um método científico rigoroso para coleta de dados
        \end{itemize}
        \item Não Probabilística
        \begin{itemize}
            \item Depende da escolha deliberada dos elementos pelo pesquisador, 
            baseada em critérios e julgamento, sem garantia de representatividade 
            estatística
        \end{itemize}
    \end{itemize}

\subsubsection{Métodos de Amostragem}
\begin{itemize}
    \item Aleatória Simples
    \begin{itemize}
        \item Definição: Todos os elementos da população possuem a mesma 
        probabilidade de serem selecionados
        \item Como realizar:
        \begin{itemize}
            \item Numere todos os elementos da população
            \[U=\{1,2,3,\dots,N\}\]
            \begin{scriptsize}

                $N$: Número de unidade amostral da população \\
                $n$: Número de unidade amostral de amostras
            \end{scriptsize}
        \end{itemize}
        \item O sorteio pode ser:
        \begin{itemize}
            \item Com reposição: o elemento retorna à população e pode ser 
            sorteado novamente
            \begin{itemize}
                \item Número de amostras possíveis: $N^n$, logo:
                \[
                P(U = u_i) = \frac{1}{N^n}
                \]
            \end{itemize}
            \item Sem reposição: o elemento não retorna e não poderá ser 
            selecionado novamente
            \begin{itemize}
                \item Número de possiveis amostras:
                \[
                C(N, n) = \left(\begin{array}{c} N \\ n \end{array}\right)
                = \frac{N!}{n!(N-n)!}
                \]
                logo,
                \[
                P(U = u_i) = \frac{1}{\left(\begin{array}{c} N \\ n \end{array}\right)}
                \]
                \item Observação: Uma combinação é uma seleção não ordenada de 
                elementos
            \end{itemize}
        \end{itemize}
        \item Vantagens: 
        \begin{itemize}
            \item Simples de aplicar
            \item Permite medir a precisão das estimativas
            \item Propriedades estatísticas bem conhecidas
        \end{itemize}
        \item Desvantagens:
        \begin{itemize}
            \item Requer lista completa da população
            \item Pode ser caro se a amostra for muito dispersa
            \item Não utiliza informações auxiliares
        \end{itemize}
        \item Qual deve ser o tamanho da amostra?
        \begin{itemize}
            \item O tamanho da amostra da
        \end{itemize}
    \end{itemize}
    \item Sistemática
    \begin{itemize}
        \item Definição: Os elementos da população estão ordenados, e a seleção 
        da amostra é realizada em intervalos regulares
        \item Etapas:
        \begin{enumerate}
            \item Defina o tamanho da população ($N$)
            \item Defina o tamanho da amostra ($n$)
            \item Calcule o intervalo da seleção ($IS$)
                \[IS = \frac{N}{n}\]
            \item Escolha aleatoriamente o primeiro elemento entre 1 e $IS$
            \item Selecione os demais elementos adicionado sucessivamente o 
            intervalo $IS$
        \end{enumerate}
    \end{itemize}
    \item Estratificada
    \begin{itemize}
        \item Definição: Divisão da população $U$ em $H$ grupos disjuntos e 
        homogêneos, conhecidos como estratos. Entre os estratos, existem 
        diferenças significativas, evidenciando a heterogeneidade da população
        \item Passos do procedimento:
        \begin{enumerate}
            \item Selecionar uma amostra dentro de cada estrato, de forma 
            independente
            \item Formar a amostra final unindo todas as amostras selecionadas 
            nos estratos
            \item Dentro de cada estrato, podem ser usados diferentes métodos 
            de coleta ou amostragem, conforme necessário
        \end{enumerate}
    \end{itemize}

\end{itemize}
    
\end{document}
